% Part: first-order-logic
% Chapter: sequent-calculus
% Section: provability-quantifiers

\documentclass[../../../include/open-logic-section]{subfiles}

\begin{document}

\olfileid[fr]{fol}{seq}{qpr}
\olsection{\usetoken{S}{derivability} et quantificateurs}

\begin{explain}
Pour démontrer le théorème de complétude, nous aurons aussi besoin
des propriétés de $\Proves$ établies dans cette section à l'aide
des règles du calcul des séquents.
\end{explain}

\begin{thm}
\ollabel{thm:strong-generalization}
Si $c$ est une constante qui ne figure ni dans $\Gamma$ ni dans
$!A(x)$, et si $\Gamma \Proves !A(c)$, alors
$\Gamma \Proves \lforall[x][!A(x)]$.
\end{thm}

\begin{proof}
Soit $\pi_0$ une !!{derivation}, dans $\Log{LK}$, de
$\Gamma_0 \Sequent !A(c)$, où $\Gamma_0 \subseteq \Gamma$
est une partie finie. En ajoutant une inférence $\RightR{\lforall}$,
on obtient une !!{derivation} de
$\Gamma_0 \Sequent \lforall[x][!A(x)]$ : la constante $c$ ne
figure ni dans $\Gamma$ ni dans $!A(x)$, de sorte que la condition
sur l'eigenvariable est satisfaite.
\end{proof}

\begin{prop}
\ollabel{prop:provability-quantifiers}
Soit $t$ un terme clos.\footnote{Note de l'édition française :
la convention sur les termes des règles de quantification est
rappelée explicitement ; elle reste implicite ici dans la source.}
\begin{tagenumerate}{prvEx,prvAll}
\tagitem{prvEx}{$!A(t) \Proves \lexists[x][!A(x)]$.}{}
\tagitem{prvAll}{$\lforall[x][!A(x)] \Proves !A(t)$.}{}
\end{tagenumerate}
\end{prop}

\begin{proof}
\begin{tagenumerate}{prvEx,prvAll}
\tagitem{prvEx}{Le séquent $!A(t) \Sequent \lexists[x][!A(x)]$
  est !!{derivable} :
  \begin{prooftree}
    \Axiom$!A(t) \fCenter !A(t)$
    \RightLabel{\RightR{\lexists}}
    \UnaryInf$!A(t) \fCenter \lexists[x][!A(x)]$
    \end{prooftree}}{}
\tagitem{prvAll}{Le séquent $\lforall[x][!A(x)] \Sequent !A(t)$
  est !!{derivable} :
  \begin{prooftree}
    \Axiom$!A(t) \fCenter !A(t)$
    \RightLabel{\LeftR{\lforall}}
    \UnaryInf$\lforall[x][!A(x)] \fCenter !A(t)$
    \end{prooftree}}{}
\end{tagenumerate}
\end{proof}

\end{document}
